% Research section prepared against ProveIt 7bd45777a8a127e5dc69aff8887ca36a79a3059d.
% This fragment uses only standard amsmath/amsthm commands.  Labels begin q4:.
% Suggested source citations are recorded in SOURCE_NOTES.md beside this file.

\section{The quartic digamma finite part and a depth-three harmonic coefficient}
\label{q4:section}

The cubic harmonic report~\cite{IncomingCubic} leaves a specific
next question: carry out the
power-to-harmonic reduction for the fourth power of the digamma function,
with the entire remainder and every endpoint constant fixed.  The theorem
below completes that reduction.  The retained depth-three coefficient is
defined explicitly; a finite evaluation of it in ordinary zeta and
Stieltjes constants is a separate arithmetic question.

Throughout this section, write $\gamma=\gamma_0$, and use
\[
 \zeta(1+u)=\frac1u+\sum_{j\geq0}\frac{(-1)^j\gamma_j}{j!}u^j.
\]
For strict multiple zeta functions the larger index is first:
\[
 Z_r(s_1,\ldots,s_r)
 =\sum_{n_1>\cdots>n_r\geq1}n_1^{-s_1}\cdots n_r^{-s_r}.
\]
Only a single spectral order varies in the following definitions:
\begin{align}
 Z_2(1+u,1)
 &=u^{-2}+\gamma u^{-1}+\frac{\gamma^2-\zeta(2)}2
       +\eta u+O(u^2),\label{q4:eta}\\
 Z_3(1+u,1,1)
 &=u^{-3}+\gamma u^{-2}
       +\frac{\gamma^2-\zeta(2)}{2u}
       +\frac{\gamma^3-3\gamma\zeta(2)+2\zeta(3)}6
       +\theta u+O(u^2).\label{q4:theta}
\end{align}
Thus $\theta$ is the linear Laurent coefficient, with no factorial or sign
included.  These conventions agree with the incoming ordered-Hurwitz and
cubic harmonic reports~\cite{IncomingOrdered,IncomingCubic}.
The expansions can also be established directly
from the finite elementary symmetric sums used below.

Define $\operatorname{FP}\int_0^1$ as the constant coefficient of
$\int_\varepsilon^1$ in the coordinate $\varepsilon$ itself.  In particular,
$\operatorname{FP}\int_0^1x^{-j}\,dx=-1/(j-1)$ for $j\geq2$, whereas
$\operatorname{FP}\int_0^1x^{-1}\,dx=0$.

\begin{theorem}[Quartic harmonic bridge]\label{q4:main}
Let
\[
 A_4=12\zeta(3)-12\gamma\zeta(2)+18\zeta(4)+4\zeta'(3).
\]
Then
\begin{equation}\label{q4:bridge}
 \boxed{\displaystyle
 Q_4:=\operatorname{FP}\int_0^1\psi(x)^4\,dx
 =A_4-24\theta+24\gamma\eta
        +12\bigl(\gamma^2-\zeta(2)\bigr)\gamma_1.}
\end{equation}
Equivalently, an absolutely convergent arithmetic representation is
\begin{equation}\label{q4:arithmetic}
 \boxed{\begin{aligned}
 Q_4={}&A_4+12\gamma_3-24\zeta(2)\gamma_1\\
 &+12\sum_{n=1}^{\infty}
 \frac{\bigl(\psi(n)^2+\psi'(n)-\log^2n\bigr)\log n}{n}.
 \end{aligned}}
\end{equation}
In generalized harmonic notation the numerator inside the parentheses is
\[
 (H_{n-1}-\gamma)^2+\zeta(2)-H_{n-1}^{(2)}-\log^2n.
\]
Using the already proved cubic bridge, one may isolate the single
depth-three coefficient in
\begin{align}\label{q4:single-coordinate}
 \operatorname{FP}\int_0^1\bigl\{\psi(x)^4+4\gamma\psi(x)^3\bigr\}\,dx
 ={}&12\zeta(3)+18\zeta(4)+4\zeta'(3)\notag\\
 &-24\theta-12\bigl(\gamma^2+\zeta(2)\bigr)\gamma_1.
\end{align}
No arithmetic independence or ordinary-constant evaluation of $\theta$ is
asserted by these formulas.
\end{theorem}

\subsection{A global quartic partial-fraction identity}

For $n\geq0$ put
\begin{equation}\label{q4:coefficients}
 \begin{split}
 b_n&=H_n-\gamma,\qquad
 c_n=\zeta(2)+H_n^{(2)},\qquad
 d_n=H_n^{(3)}-\zeta(3),\\
 q_n&=6b_n^2-4c_n,\qquad
 a_n=-b_n^3+3b_nc_n-d_n,
 \end{split}
\end{equation}
with $H_0^{(r)}=0$.  The quartic principal part at $-n$ is
\begin{equation}\label{q4:Pn}
 P_n(z)=\frac1{(z+n)^4}-\frac{4b_n}{(z+n)^3}
             +\frac{q_n}{(z+n)^2}+\frac{4a_n}{z+n}.
\end{equation}
The constant coefficient of $\psi(z)^4$ at $z=0$ is
\begin{equation}\label{q4:kappa}
 \kappa_4=\gamma^4-12\gamma^2\zeta(2)+6\zeta(2)^2
                        +12\gamma\zeta(3)-4\zeta(4).
\end{equation}

\begin{lemma}[The entire remainder is fixed]\label{q4:ML}
For $z\notin\{0,-1,-2,\ldots\}$,
\begin{equation}\label{q4:MLformula}
 \psi(z)^4=P_0(z)+\kappa_4+
                   \sum_{n=1}^{\infty}\bigl(P_n(z)-P_n(0)\bigr).
\end{equation}
The series converges normally on every compact set avoiding the poles.
\end{lemma}

\begin{proof}
The digamma recurrence and the Taylor expansion at zero
\cite[Sections~5.5, 5.7]{DLMF} give
\[
 \psi(-n+\varepsilon)=-\varepsilon^{-1}+b_n+c_n\varepsilon
           +d_n\varepsilon^2
           +\bigl(\zeta(4)+H_n^{(4)}\bigr)\varepsilon^3
           +O(\varepsilon^4).
\]
Taking its fourth power proves \eqref{q4:Pn}; its constant term at zero
is \eqref{q4:kappa}.  Matching principal parts by itself would leave an
unknown entire function, so one additional argument is necessary.

Let $f(z)=\psi(z)^4-P_0(z)$, which is regular at zero.  For fixed $z$,
integrate
\[
 \frac{z f(w)}{w(w-z)}
\]
over $|w|=N+1/2$, with $N$ sufficiently large.  On these circles,
$f(w)=O(\log^4N)$ uniformly.  In the right half-plane this follows from
the usual digamma asymptotic~\cite[Section~5.11]{DLMF}.
In the left half-plane use
$\psi(w)=\psi(1-w)-\pi\cot(\pi w)$.  The cotangent is bounded uniformly
on the chosen circles: for $|\operatorname{Im}w|\geq1/4$ its exponential
formula suffices, and for the remaining arcs the real part remains a
fixed distance from every integer.  Hence the contour integral is
$O(\log^4N/N)$ and tends to zero.

Its residues at $z$ and $0$ are $f(z)$ and $-f(0)$, and the residue at
$-n$ is $-(P_n(z)-P_n(0))$.  The residue theorem proves
\eqref{q4:MLformula}.  Finally,
$b_n=O(\log n)$, $q_n=O(\log^2n)$ and $a_n=O(\log^3n)$, while
$(n+z)^{-1}-n^{-1}=O(n^{-2})$ uniformly on a fixed compact set.
This proves normal convergence as well as the identity.
\end{proof}

\subsection{Integration and finite summation by parts}

Integrating the normally convergent remainder in \eqref{q4:MLformula}
term by term gives
\begin{align}\label{q4:integrated}
 Q_4={}&-\frac13-2\gamma-q_0+\kappa_4\notag\\
 &+\sum_{n\geq1}\left\{
   \frac{n^{-3}-(n+1)^{-3}}3-n^{-4}
   -2b_n\bigl(n^{-2}-(n+1)^{-2}\bigr)+4b_nn^{-3}\right.\notag\\
 &\hspace{34mm}\left.
   +q_n\bigl(n^{-1}-(n+1)^{-1}-n^{-2}\bigr)
   +4a_n\left(\log\left(1+\frac1n\right)-\frac1n\right)
  \right\}.
\end{align}
The finite parts $-1/3$, $-2\gamma$ and $-q_0$ are retained explicitly.
The simple-pole finite part at the origin is zero.

We use the following classical convergent Euler sums
\cite[Equations~25.16.8 and 25.16.13]{DLMF}:
\begin{align*}
 \sum_{n\geq1}\frac{H_n}{n^2}&=2\zeta(3),&
 \sum_{n\geq1}\frac{H_n}{n^3}&=\frac54\zeta(4),\\
 \sum_{n\geq1}\frac{H_n^2}{n^2}&=\frac{17}4\zeta(4),&
 \sum_{n\geq1}\frac{H_n^{(2)}}{n^2}
     &=\frac{\zeta(2)^2+\zeta(4)}2.
\end{align*}
The first three can be taken from the standard Euler-sum formulas; the
fourth follows by splitting the product $\zeta(2)^2$ into its two strict
triangles and its diagonal.  Also $\zeta(2)^2=5\zeta(4)/2$.
The exact differences
\begin{equation}\label{q4:differences}
 b_n-b_{n-1}=\frac1n,\qquad
 q_n-q_{n-1}=\frac{12b_{n-1}}n+\frac2{n^2},\qquad
 a_n-a_{n-1}=\frac{3(c_{n-1}-b_{n-1}^2)}n+\frac1{n^3}
\end{equation}
then reduce all terms in \eqref{q4:integrated} except the logarithmic
ones to
\begin{equation}\label{q4:base}
 Q_4=B_4+4\sum_{n\geq1}a_n
          \left(\log\left(1+\frac1n\right)-\frac1n\right),
\end{equation}
where
\[
 B_4=\gamma^4-18\gamma^2\zeta(2)-12\gamma\zeta(2)
             +32\gamma\zeta(3)+12\zeta(3)+\frac{13}2\zeta(4).
\]
For a directly checkable bookkeeping formula, the lower-pole terms before
substitution of the Euler sums are
\[
 \kappa_4-\zeta(4)
 +4\sum_{n\geq1}\frac{b_n}{n^3}
 +12\sum_{n\geq1}\frac{b_{n-1}}{n^2}
 -6\sum_{n\geq1}\frac{b_n^2}{n^2}
 +4\sum_{n\geq1}\frac{c_n}{n^2}.
\]
This equals $B_4$ exactly.

The residue difference in \eqref{q4:differences} is the decisive step.
For a finite integer $N\geq1$, define
\[
 M_N=\sum_{n=1}^N
             \frac{(b_{n-1}^2-c_{n-1})\log n}{n}.
\]
Finite summation by parts gives the exact equality
\begin{equation}\label{q4:logsum}
 \sum_{n=1}^N a_n\log\left(1+\frac1n\right)
       =a_N\log(N+1)+3M_N-\sum_{n=1}^N\frac{\log n}{n^3}.
\end{equation}
To evaluate the other finite sum, let $e_j(N)$ be the elementary symmetric
sum of degree $j$ in $1,1/2,\ldots,1/N$, and write
\[
 H_N(r_1,\ldots,r_d)
  =\sum_{N\geq n_1>\cdots>n_d\geq1}
         n_1^{-r_1}\cdots n_d^{-r_d}.
\]
The polynomial expression for the residue is
\begin{equation}\label{q4:elementary}
 a_n=-6e_3(n)+6\gamma e_2(n)+3(\zeta(2)-\gamma^2)H_n
           +\gamma^3-3\gamma\zeta(2)+\zeta(3)+H_n^{(3)}.
\end{equation}
Use $e_j(n)=e_j(n-1)+e_{j-1}(n-1)/n$ in each sum, and split the
product $H_NH_N^{(3)}$ into strict triangles.  This proves, without
any infinite rearrangement,
\begin{align}\label{q4:finitesum}
 \sum_{n=1}^N\frac{a_n}{n}
 ={}&-6e_4(N)+6\gamma e_3(N)+3(\zeta(2)-\gamma^2)e_2(N)\notag\\
 &+(\gamma^3-3\gamma\zeta(2)+\zeta(3))H_N+H_NH_N^{(3)}\notag\\
 &-6H_N(2,1,1)+6\gamma H_N(2,1)
       +3(\zeta(2)-\gamma^2)H_N^{(2)}-H_N(3,1).
\end{align}
The only multiple-zeta limits needed here are
\[
 Z_3(2,1,1)=\zeta(4),\qquad
 Z_2(2,1)=\zeta(3),\qquad
 Z_2(3,1)=\frac14\zeta(4).
\]
For example,
\[
 Z_3(2,1,1)=\frac16\int_0^1\frac{[-\log(1-t)]^3}{t}\,dt
           =\frac16\int_0^1\frac{(-\log t)^3}{1-t}\,dt
           =\zeta(4),
\]
and the double-zeta values follow from the displayed Euler sums.

The elementary symmetric formulas are
\[
 \begin{split}
 e_2(N)&=\frac{H_N^2-H_N^{(2)}}2,\\
 e_3(N)&=\frac{H_N^3-3H_NH_N^{(2)}+2H_N^{(3)}}6,\\
 e_4(N)&=\frac{H_N^4-6H_N^2H_N^{(2)}+3(H_N^{(2)})^2
                +8H_NH_N^{(3)}-6H_N^{(4)}}{24}.
 \end{split}
\]
Substitute these into \eqref{q4:finitesum}, using
$H_N=\log N+\gamma+O(N^{-1})$ and
$H_N^{(r)}=\zeta(r)+O(N^{1-r})$ for $r\geq2$.
All discarded terms are $o(1)$, including their polynomial logarithmic
factors.  The result, including its full constant term, is
\begin{equation}\label{q4:asymptoticsum}
 \sum_{n=1}^N\frac{a_n}{n}
 =-\frac{\log^4N}{4}+3\zeta(2)\log^2N+C_4+o(1),
\end{equation}
where
\begin{equation}\label{q4:C4}
 C_4=\frac{\gamma^4}{4}-\frac92\gamma^2\zeta(2)
                    +8\gamma\zeta(3)-\frac{23}8\zeta(4).
\end{equation}
Also
\[
 a_N\log(N+1)=-\log^4N+6\zeta(2)\log^2N+o(1).
\]
Equations \eqref{q4:base}--\eqref{q4:asymptoticsum} now yield
\begin{equation}\label{q4:limit}
 Q_4=A_4+12\lim_{N\to\infty}
       \left\{M_N-\frac{\log^4N}{4}+\zeta(2)\log^2N\right\}.
\end{equation}
Indeed $B_4-4C_4=12\zeta(3)-12\gamma\zeta(2)+18\zeta(4)$,
and $-\sum\log n/n^3=\zeta'(3)$.  This is where an omitted endpoint
constant or a sign error in the logarithmic sum would alter the theorem.

\subsection{Identification of the Laurent coefficient}

Since
\[
 b_{n-1}^2-c_{n-1}
   =\psi(n)^2+\psi'(n)-2\zeta(2),
\]
and $\psi(n)=\log n+O(n^{-1})$, $\psi'(n)=O(n^{-1})$, the series
in \eqref{q4:arithmetic} is absolutely convergent, with summands
$O(\log^2n/n^2)$.  The standard Stieltjes limits give
\begin{align}\label{q4:Mconstant}
 \lim_{N\to\infty}
       \left\{M_N-\frac{\log^4N}{4}+\zeta(2)\log^2N\right\}
  ={}&\gamma_3-2\zeta(2)\gamma_1\notag\\
 &+\sum_{n\geq1}
       \frac{(\psi(n)^2+\psi'(n)-\log^2n)\log n}{n}.
\end{align}
Together with \eqref{q4:limit}, this proves \eqref{q4:arithmetic}.

For the precise spectral interpretation, define the holomorphic function
\[
 D_4(s)=\sum_{n\geq1}
       \frac{\psi(n)^2+\psi'(n)-\log^2n}{n^s},
       \qquad\operatorname{Re}s>0.
\]
The series and every fixed derivative converge normally on compact
subsets of this half-plane.  In $\operatorname{Re}s>1$ the finite identity
$H_{n-1}^2-H_{n-1}^{(2)}=2e_2(n-1)$ gives
\begin{equation}\label{q4:Didentity}
 D_4(s)=2Z_3(s,1,1)-2\gamma Z_2(s,1)
                  +(\gamma^2+\zeta(2))\zeta(s)-\zeta''(s).
\end{equation}
Analytic continuation carries the identity to the germ at $s=1$.
Comparing its linear coefficients, with \eqref{q4:eta} and
\eqref{q4:theta}, yields
\[
 D_4'(1)=2\theta-2\gamma\eta
                    -(\gamma^2+\zeta(2))\gamma_1+\gamma_3.
\]
The convergent sum in \eqref{q4:arithmetic} is $-D_4'(1)$.
Substitution proves \eqref{q4:bridge}.  Finally, the incoming cubic
report~\cite{IncomingCubic}
proves
\[
 Q_3=\operatorname{FP}\int_0^1\psi(x)^3\,dx
    =3\zeta(2)-6\eta-6\gamma\gamma_1.
\]
Adding $4\gamma Q_3$ gives \eqref{q4:single-coordinate}.

\paragraph{A direct existence check for the depth-three Laurent germ.}
There is also a useful elementary Mellin representation:
\begin{equation}\label{q4:heightone-mellin}
 Z_3(s,1,1)=\frac1{2\Gamma(s)}\int_0^\infty
     \frac{t^{s-1}\log^2(1-e^{-t})}{e^t-1}\,dt,
     \qquad\operatorname{Re}s>1.
\end{equation}
Indeed the generating series of $e_2(n-1)$ is
$\log^2(1-z)/(2(1/z-1))$.  Mellin integration is absolutely convergent
in the stated domain.  Near zero the kernel is
$\log^2t/(2t)+O(\log^2t)$, so subtracting that one local term gives
a holomorphic remainder $R_3(u)$ in
\[
 \Gamma(1+u)Z_3(1+u,1,1)=u^{-3}+R_3(u).
\]
The identity
$\frac{d}{dt}\log^3(1-e^{-t})=3\log^2(1-e^{-t})/(e^t-1)$
shows $R_3(0)=0$.  The inverse-Gamma expansion therefore gives exactly
the principal and constant coefficients in \eqref{q4:theta}.  In
particular,
\begin{align}\label{q4:theta-kernel}
 \theta={}&\frac{\gamma^4-6\gamma^2\zeta(2)+3\zeta(2)^2
                       +8\gamma\zeta(3)-6\zeta(4)}{24}\notag\\
 &+\frac12\int_0^1\log t
       \left\{\frac{\log^2(1-e^{-t})}{e^t-1}
                        -\frac{\log^2t}{t}\right\}\,dt\notag\\
 &+\frac12\int_1^\infty
       \frac{\log t\,\log^2(1-e^{-t})}{e^t-1}\,dt.
\end{align}
Both integrals in this display converge ordinarily.  This fixes the
retained coordinate by an elementary, one-dimensional kernel as well as
by the strict Laurent definition; it does not remove its arithmetic
content by renaming it.

\subsection{A normalized quartic antiderivative and derivative identities}

For $j\geq2$, define
\[
 F_j(n,z)=\frac{(n+z)^{1-j}-n^{1-j}}{1-j}-\frac{z}{n^j},
 \qquad
 F_1(n,z)=\log\left(1+\frac zn\right)-\frac zn.
\]
Use the logarithm real on the positive real axis, continued to
$\operatorname{Re}z>0$.  Put
$p_{4,n}=1$, $p_{3,n}=-4b_n$, $p_{2,n}=q_n$, $p_{1,n}=4a_n$.

\begin{corollary}[Normalized primitive]\label{q4:primitive}
On $\operatorname{Re}z>0$, the function
\begin{align}\label{q4:primitiveformula}
 \mathcal P_4(z)={}&-\frac1{3z^3}-\frac{2\gamma}{z^2}
       -\frac{6\gamma^2-4\zeta(2)}z
       +4(\gamma^3-3\gamma\zeta(2)+\zeta(3))\log z
       +\kappa_4z\notag\\
 &+\sum_{n=1}^\infty\sum_{j=1}^4p_{j,n}F_j(n,z)
\end{align}
satisfies $\mathcal P_4'(z)=\psi(z)^4$.  Its expansion at the singular
origin has constant term zero.  Consequently, for every real $b>0$,
\[
 \operatorname{FP}\int_0^b\psi(x)^4\,dx=\mathcal P_4(b),
 \qquad \mathcal P_4(1)=Q_4,
\]
and, for $0<a<b$, the ordinary integral is
$\int_a^b\psi(x)^4\,dx=\mathcal P_4(b)-\mathcal P_4(a)$.
\end{corollary}

\begin{proof}
Each $F_j$ is the primitive vanishing at zero of
$(n+z)^{-j}-n^{-j}$.  Uniformly on a fixed compact set,
$F_j(n,z)=O(n^{-j-1})$.  The coefficient growth established in
Lemma~\ref{q4:ML} therefore proves normal convergence of the series and
its differentiated series.  Differentiation gives
\eqref{q4:MLformula}.  The same estimates give uniform convergence on
each disk $|z|\leq r<1$, including the origin, where every summand
vanishes.  Thus the displayed singular terms fix the constant
coefficient to be zero.
\end{proof}

\begin{corollary}[Quartic polygamma hierarchy]\label{q4:derivatives}
For every integer $r\geq1$, away from the poles of $\psi$,
\begin{align}\label{q4:allpolygamma}
 &\sum_{r_1+r_2+r_3+r_4=r}
       \frac{r!}{r_1!r_2!r_3!r_4!}
       \prod_{j=1}^4\psi^{(r_j)}(z)\notag\\
 &\quad=(-1)^rr!\sum_{n=0}^{\infty}
       \left\{\frac{\binom{r+3}{3}}{(n+z)^{r+4}}
       -\frac{4b_n\binom{r+2}{2}}{(n+z)^{r+3}}
       +\frac{(r+1)q_n}{(n+z)^{r+2}}
       +\frac{4a_n}{(n+z)^{r+1}}\right\}.
\end{align}
The series converges normally.  At $z=1$ its value is a finite polynomial
in $\gamma,\zeta(2),\ldots,\zeta(r+1)$.
\end{corollary}
\begin{proof}
Differentiate \eqref{q4:MLformula} $r$ times.  Its constant subtractions
disappear, and the resulting terms are bounded by a constant times
$n^{-r-1}(1+\log^3n)$ on each compact set avoiding the poles.
The product rule gives the left side.  The last assertion follows from
$\psi(1)=-\gamma$ and
$\psi^{(k)}(1)=(-1)^{k+1}k!\zeta(k+1)$ for $k\geq1$.
\end{proof}

\subsection{Independent verification and a reproducible tail formula}

The supplied program checks 44 exact algebraic identities, including
all principal coefficients, the residue differences, the finite
summation identity for $1\leq N\leq25$, the endpoint-constant
cancellation, Laurent coefficient extraction, and the primitive
derivatives.  Its separate numerical calculations use 100 decimal
digits and give
\[
 Q_4=13.510120376015827411606240712830448441932798533301918600986096436\ldots.
\]
One evaluator uses the arithmetic series with a Binet tail.  The other
integrates the convergent Taylor series of $(x\psi(x))^4$ near zero and
uses ordinary quadrature away from zero; it does not use the harmonic
bridge or the Mittag--Leffler expansion.  Two choices of the matching
point, $1/8$ and $1/16$, agree to the displayed precision.  The finest
arithmetic result differs from these evaluations by less than
$1.40\cdot10^{-83}$.

Here is the explicit analytic tail control used in the arithmetic
calculation.  For $K\geq0$ set
\[
 a_K(x)=-\frac1{2x}-\sum_{k=1}^K\frac{B_{2k}}{2k\,x^{2k}},
 \qquad
 p_K(x)=\frac1x+\frac1{2x^2}
                      +\sum_{k=1}^K\frac{B_{2k}}{x^{2k+1}}.
\]
Binet's integral~\cite[Equation~5.9.15]{DLMF} and the finite
geometric expansion give, for $x>0$,
\[
 |\psi(x)-\log x-a_K(x)|\leq E_Kx^{-2K-2},\qquad
 |\psi'(x)-p_K(x)|\leq G_Kx^{-2K-3},
\]
where $G_K=|B_{2K+2}|$ and $E_K=G_K/(2K+2)$.
For the differentiated bound, expand $(1+y)^{-2}$ under the differentiated
Binet integral; its Taylor remainder after $K$ terms has magnitude at
most $(K+1)y^K$ for $y\geq0$.

For $N\geq2$ put
$U_{N,K}=1/(2N)+\sum_{k=1}^K|B_{2k}|/(2kN^{2k})$.
Replace $\psi(n)^2+\psi'(n)-\log^2n$ in the tail $n\geq N$ by
\[
 2\log n\,a_K(n)+a_K(n)^2+p_K(n).
\]
This finite polynomial in $n^{-1}$ and $\log n$ sums exactly by the first
two spectral derivatives of $\zeta(s,N)$.  The error in the series
itself is bounded by
\begin{align}\label{q4:tailbound}
 &2E_K\zeta''(2K+3,N)
    -2U_{N,K}E_K\zeta'(2K+3,N)\notag\\
 &\qquad{}-E_K^2\zeta'(4K+5,N)
    -G_K\zeta'(2K+4,N).
\end{align}
To obtain the error in $Q_4$, multiply this bound by $12$.
Indeed, if $R=\psi-\log-a_K$, then
\[
 |\psi^2+\psi'-(\log+a_K)^2-p_K|
 \leq2(\log n+U_{N,K})E_Kn^{-2K-2}
          +E_K^2n^{-4K-4}+G_Kn^{-2K-3},
\]
and summing after multiplication by $\log n/n$ gives
\eqref{q4:tailbound}.  With $N=64$, $K=34$, the bound for $Q_4$ is
below $1.88\cdot10^{-83}$.  The implementation uses ordinary mpmath
arithmetic, so this analytic truncation bound does not include an
interval enclosure of rounding error; the numerical evidence is
reported with that qualification.

\paragraph{Scope of the advance.}
The Mittag--Leffler method, harmonic Laurent coefficients, the Mellin
representation, and the Euler sums are classical.  The incoming cubic
report supplies the immediate proof strategy and explicitly asks for
the quartic reduction.  The new result relative to the inspected corpus
is \eqref{q4:bridge}, including its exact finite correction, convergent
arithmetic form and normalized primitive.  The claim is a solution of
that stated quartic reduction problem, not a proof that $\theta$ is
irreducible or a claim of global priority for every equivalent formula.

\paragraph{Further questions.}
The contour argument extends to each fixed power $\psi^k$, but the finite
arithmetic reduction has only been carried out here through $k=4$.
The next case asks which strict harmonic Laurent coefficients occur in
$\operatorname{FP}\int_0^1\psi(x)^5\,dx$ and whether some canonical
polynomial in $\psi$ cancels all lower-depth coefficients.  A second
question is to combine the new quartic bridge with the cubic logarithmic
sine moment supplied by the existing beta generator and determine the
smallest remaining spectral coordinate space.  A third is to obtain a
common-Hurwitz-shift version with all boundary terms fixed; simply
replacing $\gamma_j$ by $\gamma_j(a)$ in \eqref{q4:bridge} is not
justified by this proof.
